% Appendix A: full proofs (r470). Included from paper.tex.
\section{Proofs}
\label{app:proofs}

\subsection{Monotonicity of the fixed-$k$ flip probability (Lemma~\ref{lem:exact})}

Recall
\[
f_K(k)=\P\big(\MV(k)\neq\MV(N)\,\big|\,K\big),
\qquad x_k\sim\HG(N,K,k),
\]
with the negative side winning ties ($\MV(k)=1$ iff $x_k>k/2$).
We prove that $f_K(k)$ is nonincreasing along the odd numbers
$k\in\{1,3,5,\dots,N-1\}$ for every fixed $K$; this suffices for the paper,
since the fixed-$k$ certificate family is evaluated on odd budgets only, and
it makes $k^*(\alpha)=\min\{k:\E_H[f_K(k)]\le\alpha\}$ well defined and
binary-searchable for every mixture $H$ (mixtures of nonincreasing functions
are nonincreasing).

\paragraph{Setup and case split.}
$\MV(k)\neq\MV(N)$ can happen in only two ways:
\begin{itemize}
\item \textbf{Case U (upset).} $K\le N/2$ (full vote negative) and $x_k>k/2$
(prefix votes positive).
\item \textbf{Case R (reversal).} $K>N/2$ (full vote positive) and $x_k\le k/2$
(prefix votes negative).
\end{itemize}
Write $f_K(k)=U_K(k)+R_K(k)$ with $U_K(k)=\P(x_k>k/2)\,\mathbf 1\{K\le N/2\}$
and $R_K(k)=\P(x_k\le k/2)\,\mathbf 1\{K>N/2\}$. At most one of the two terms
is nonzero for any $K$.

\paragraph{The reflection identity.}
Swapping the labels ``pass''/``fail'' of all $N$ rollouts maps the law
$\HG(N,K,k)$ of $x_k$ onto the law of $k-x_k$ under count $N-K$, i.e.\
$\P_{K}(x_k>s)=\P_{N-K}(x_k<k-s)$. Hence, for odd $k$,
\[
U_K(k)=\P_K(x_k>k/2)=\P_{N-K}(x_k<k/2)=R_{N-K}(k)
\qquad(K\le N/2),
\]
since $k$ odd makes $k/2$ a half-integer, so the strict and weak
inequalities coincide. Every upset probability thus equals a reversal
probability at the reflected count, and it suffices to prove
\begin{equation}
\label{eq:revmain}
R_K(k+2)\le R_K(k)\qquad\text{for all }K>N/2\text{ and odd }k,
\end{equation}
which covers Case R directly and Case U by reflection (the same display
applied at count $N-K$).

\paragraph{Coupling construction.}
Fix $K>N/2$ and odd $k\le N-3$. We couple the two hypergeometric draws on a
single probability space. Let the population of $N$ rollouts consist of $K$
passes and $N-K$ fails. Draw an \emph{ordered} sample of size $k+2$ uniformly
without replacement and split it as
\[
\underbrace{X_1,\dots,X_k}_{\text{prefix }k}\;\Big|\;
\underbrace{X_{k+1},X_{k+2}}_{\text{extra pair }(A,B)} .
\]
The marginal law of $(X_1,\dots,X_k)$ is exactly $\HG(N,K,k)$, and
$x_{k+2}=x_k+A+B$ has law $\HG(N,K,k+2)$. All statements below are about this
joint space, so they hold for the two marginals simultaneously.

Let $t=k/2+1/2$ be the (integer) flip threshold: the prefix at $k$ votes
negative iff $x_k\le t-1$; the prefix at $k+2$ votes negative iff
$x_{k+2}\le t$. Partition the event $\{x_k=s\}$ by the value of the pair
$A+B\in\{0,1,2\}$:
\[
\begin{aligned}
E_0&=\{x_k=t-1,\ A+B=0\}, &&\text{flip at }k,\ \text{flip at }k+2,\\
E_1&=\{x_k=t-1,\ A+B=1\}, &&\text{flip at }k,\ \text{flip at }k+2,\\
E_2&=\{x_k=t-1,\ A+B=2\}, &&\text{flip at }k,\ \text{no flip at }k+2,\\
F_0&=\{x_k=t,\ \ \ A+B=0\}, &&\text{no flip at }k,\ \text{flip at }k+2,\\
F_1&=\{x_k=t,\ \ \ A+B=1\}, &&\text{no flip at }k,\ \text{no flip at }k+2.
\end{aligned}
\]
All other states agree on both sides (e.g.\ $x_k\le t-2$ flips under both, since
$A+B\le2$; $x_k\ge t+1$ flips under neither). Therefore
\begin{equation}
\label{eq:rdiff}
R_K(k)-R_K(k+2)=\P(E_2)-\P(F_0),
\end{equation}
and it remains to show $\P(E_2)\ge\P(F_0)$.

\paragraph{Comparison of the two boundary atoms.}
Both atoms factor as (prefix probability) $\times$ (pair probability), and
each factor is a ratio of binomial coefficients. With $k=2t-1$ the algebra
closes completely. First,
\[
\P(x_k=t-1)=\frac{\binom{K}{t-1}\binom{N-K}{t}}{\binom{N}{k}},
\qquad
\P(x_k=t)=\frac{\binom{K}{t}\binom{N-K}{t-1}}{\binom{N}{k}},
\]
since $k-(t-1)=t$ and $k-t=t-1$. Given $x_k=t-1$, the remainder holds
$K-t+1$ passes out of $N-k$, so
$\P(A+B=2\mid x_k=t-1)=\binom{K-t+1}{2}\big/\binom{N-k}{2}$; given $x_k=t$,
the remainder holds $N-K-t+1$ fails, so
$\P(A+B=0\mid x_k=t)=\binom{N-K-t+1}{2}\big/\binom{N-k}{2}$. The
denominators $\binom{N}{k}\binom{N-k}{2}$ are common to both atoms, hence
\[
\frac{\P(E_2)}{\P(F_0)}
=\frac{\binom{K}{t-1}\binom{N-K}{t}\binom{K-t+1}{2}}
{\binom{K}{t}\binom{N-K}{t-1}\binom{N-K-t+1}{2}}
=\frac{K-t}{N-K-t},
\]
where the last equality uses the three elementary ratios
$\binom{K}{t-1}\big/\binom{K}{t}=\frac{t}{K-t+1}$,
$\binom{N-K}{t}\big/\binom{N-K}{t-1}=\frac{N-K-t+1}{t}$, and
$\binom{K-t+1}{2}\big/\binom{N-K-t+1}{2}=\frac{(K-t+1)(K-t)}{(N-K-t+1)(N-K-t)}$,
whose product telescopes to the displayed ratio. Therefore, with
$C:=\binom{K}{t}\binom{N-K}{t-1}\binom{N-K-t+1}{2}\big/
\big[\binom{N}{k}\binom{N-k}{2}\big]\ge 0$,
\begin{equation}
\label{eq:atomdiff}
\P(E_2)-\P(F_0)=C\big[(K-t)-(N-K-t)\big]=C\,(2K-N)\;\ge\;0,
\end{equation}
because $K>N/2$ in Case R (and if $t>K$ then $\P(x_k=t-1)=0$, so both atoms
vanish). By \eqref{eq:rdiff} this proves \eqref{eq:revmain} and, by the
reflection identity, $U_K(k+2)\le U_K(k)$ as well, completing the proof of
monotonicity along odd $k$.\qed

\begin{remark}[Verification at $N=32$]
Independently of the proof, the monotonicity was verified by exhaustive
enumeration over all $33\times15$ $(K,k)$ cells at $N=32$ (zero violations,
tolerance $10^{-12}$; artifact \texttt{L2\_lemmas\_r468\_result.json},
odd-$k$ grid $\{3,5,\dots,31\}$).
\end{remark}

\begin{remark}[Why only odd $k$ enters the certificate]
Requiring $k$ odd removes tie-breaking from the certificate family
(every prefix has a strict majority side), and the binary search for
$k^*(\alpha)$ is over the odd ladder only. Nothing else in the paper uses
monotonicity at even $k$.
\end{remark}

\subsection{Forced-stop bookkeeping (Theorem~\ref{thm:tower})}

The stopping rule is $\tau=\min\{k\ge 3:\cert(x_k,k)\le\alpha\}$, truncated
at $k=N$ (forced stop). The tower argument of Theorem~\ref{thm:tower} applies
to the truncated rule verbatim: $\tau$ is still a stopping time with respect
to the prefix filtration, and
\[
\P(\flip)=\E\big[\cert(x_\tau,\tau)\big]
=\E\big[\cert\,\mathbf 1\{\tau<N\}\big]+\E\big[\cert\,\mathbf 1\{\tau=N\}\big].
\]
The first term is bounded by $\alpha\,\P(\tau<N)\le\alpha$ pointwise, because
the rule stops only at states with $\cert\le\alpha$. The second term is
\emph{not} bounded by $\alpha$ pointwise: at a forced stop the prefix
certificate may exceed $\alpha$ (and can be as large as the posterior mass on
the opposite side). This is precisely why the finite-sample guarantee is
stated at the \emph{population} level (Theorem~\ref{thm:pop}): the exact
per-task flip probability $g(K;\alpha,\widehat H_m)$ computed by the DP of
Appendix~\ref{app:dp} includes the forced-stop contribution
$\E[\cert\,\mathbf 1\{\tau=N\}\mid K]$ exactly, and the empirical-Bernstein
layer then bounds the population mean $\sum_K H[K]g(K)$. Conditional
per-problem validity at forced stops is not claimed (Limitation (v)).

\begin{remark}[The forced-stop mass is small in practice]
On the TEST split, the exact fraction of paths reaching $k=N$ under
BAYES-H is $0.59\%$ at $\alpha=0.05$ and $1.10\%$ at $\alpha=0.02$ (the
certificate stops early exactly on the bimodal mass); their contribution
is included in $g$ exactly regardless (Appendix~\ref{app:dp}).
\end{remark}

\subsection{Linearity of the population flip rate}

For fixed rule and fixed prior $\widehat H_m$ used \emph{inside} the
certificate, the per-task flip probability $g(K)$ depends only on the true
count $K$ (via the hypergeometric law of prefixes) and on $\widehat H_m$
(through the posterior threshold), but is a fixed function once $\widehat H_m$
is frozen after FIT. The population replay flip rate is therefore
\[
\P_{\text{replay},H}(\flip)=\sum_{K=0}^{N} H[K]\,g(K),
\]
a \emph{linear} functional of the true mixture $H$. Consequences used in the
main text: (i) no Jensen gap --- bounding $\bar g_{\calD}$ by empirical
Bernstein bounds the population quantity directly, with estimation error of
$\widehat H_m$ entering only through the empirical mean of a bounded per-task
quantity $g\in[0,1]$; (ii) graceful degradation --- if the deployment mixture
shifts from $H$ to $H'$, $|\P_{H'}-\P_H|\le\|H'-H\|_1$ since
$g\in[0,1]$; (iii) Bonferroni over the pre-declared family of $J=64$ rules
legitimizes the min-$k$ selection, because simultaneous UCBs on all $J$
linear functionals imply the UCB of the selected rule.
